
% !TeX root = ../main.tex
\begin{chineseabstract}

随着互联网视频数据的持续增长，从长视频中自动提取语义完整且叙事连贯的关键内容已成为视频理解领域的重要研究问题。该任务面临两方面核心挑战：其一，基于底层视觉特征的无监督方法难以跨越视觉表征与高层语义之间的鸿沟，在主题刻画与叙事保持方面存在固有局限；其二，单一评分模型难以同时兼顾叙事推进、视觉显著性、情绪表达与信息密度等多维评价需求，且面对不同类型视频缺乏自适应的评估策略。此外，基于监督学习的方法对大规模人工标注依赖较强，标注本身具有较强主观性，进一步制约了模型在开放域场景中的泛化能力。针对上述问题，本文以免训练视频摘要为目标，分别从跨模态语义对齐与多维协同推理两条路径展开研究，构建基于大语言模型的视频摘要方法体系。其中，前者侧重通过跨模态语义对齐弥合视觉表征与语义理解之间的差距，后者侧重通过多角色结构化推理提升重要性评估的全面性与策略适应性。

首先，针对底层视觉特征难以有效表征视频主题语义的问题，本文提出一种基于层级语义建模的免训练视频摘要方法。该方法通过视频结构化分段、帧级视觉描述生成、场景级与视频级文本逐层聚合，将视频内容映射至统一的文本语义空间，并利用大语言模型的零样本推理能力完成全局主题提取与多粒度语义相关性建模，结合全局主题约束与局部视觉证据实现帧级重要性估计，从而将视频摘要从视觉代表性筛选转化为语义理解与匹配问题。进一步地，针对单一评分模型难以兼顾叙事推进、视觉显著性、情绪表达与信息密度等多维评价目标的问题，本文提出一种面向视频摘要的多智能体协同推理方法。该方法将单模型统一评分拆解为全局规划与多角色协作的结构化推理流程，由Planner智能体完成视频主题推断与专家权重动态分配，再由Story、Visual、Emotion与Information四类专家智能体从互补维度进行协同评分，并引入记忆机制增强时序一致性，最终经段到帧映射与归一化处理生成帧级摘要结果。

在SumMe与TVSum数据集上的实验表明，上述两种方法均可在无需额外训练的条件下取得稳定且具有竞争力的摘要性能，消融实验进一步验证了层级描述聚合、语义评分、多专家协同及后处理策略等关键模块的有效性。

\chinesekeywordstype{视频摘要；大语言模型；免训练；层级语义建模；跨模态对齐；多智能体协同推理}{应用研究}

\end{chineseabstract}

